% =============================================================================
% Appendix: Rubric-Guided Family-Specific Token Classifier
% Requires (preamble):
%   \usepackage{tcolorbox}\tcbuselibrary{breakable,skins,shadows}
%   \usepackage{amsmath,amssymb}
%   \usepackage{enumitem}
%   \usepackage{booktabs}
% Referenced from main text as Appendix~\ref{app:classifier}.
% =============================================================================

\section{Rubric-Guided Family-Specific Token Classifier}
\label{app:classifier}

This appendix details the offline construction of the family-specific token
classifier used in Section~\ref{sec:token_classification}.
Let $s\!=\!v_n^{S}$ and $t\!=\!v_n^{T}$ denote the student and teacher top-1
tokens at a high-KL position $n$ (after lightweight normalization:
lowercasing and stripping surrounding punctuation).
The classifier maps each pair $(s,t)$ to one of six categories
$\{\texttt{pivot},\texttt{intent},\texttt{risk},\texttt{function},
\texttt{consistent},\texttt{other}\}$.
At training time, risk is further split into same-surface
(\texttt{risk\_lexicon}) and disagree (\texttt{risk\_wo\_same}) so that only the
latter enters the default keep set
$\{\texttt{pivot},\texttt{intent},\texttt{risk\_wo\_same}\}$.

\paragraph{Notation.}
We write $\mathcal{L}_{\mathrm{pivot}}^{S/T}$,
$\mathcal{L}_{\mathrm{intent}}^{S/T}$,
$\mathcal{L}_{\mathrm{risk}}$, and
$\mathcal{L}_{\mathrm{func}}$ for the family-specific lexicons, and
$\mathcal{P}_{\mathrm{intent}}\subseteq
\mathcal{L}_{\mathrm{intent}}^{S}\!\times\!\mathcal{L}_{\mathrm{intent}}^{T}$
for the extracted intent pairs.
These sets are \emph{not} hand-written: they are produced by the pipeline in
Appendix~\ref{app:classifier:pipeline}.

% ---------------------------------------------------------------------------
\subsection{Universal Semantic Rubrics}
\label{app:classifier:rubrics}

We first fix a model-agnostic semantic rubric that defines what each category
\emph{means}.
An LLM-as-a-judge is prompted with this rubric (together with the user query and
student response) when labeling disagreeing $(s,t)$ pairs.
Surface-equal pairs ($s\!=\!t$) are assigned to \texttt{consistent} by a
deterministic rule and are never sent to the judge.
Pairs that match none of the four extractable categories are treated as
\texttt{other} (no lexicon is accumulated for \texttt{other}).
When categories conflict, the priority is
\texttt{pivot}$>$\texttt{risk}$>$\texttt{intent}$>$\texttt{function}.

% Card style: full-width gray title bar + soft drop shadow (cf. pivot-rules figure).
\tcbset{
  rubricbox/.style={
    enhanced,
    breakable,
    colback=white,
    colframe=black!55,
    boxrule=0.6pt,
    arc=2.5pt,
    left=7pt,
    right=7pt,
    top=5pt,
    bottom=6pt,
    toptitle=4pt,
    bottomtitle=4pt,
    lefttitle=7pt,
    righttitle=7pt,
    colbacktitle=black!10,
    coltitle=black,
    fonttitle=\bfseries\small,
    shadow={1.2mm}{-1.2mm}{0mm}{black!18},
  },
}

\begin{tcolorbox}[rubricbox,
  title={Category 1: Pivot ($\mathcal{R}_{\mathrm{pivot}}$)}]
Identifies structural turns from compliance / task planning toward refusal
or a safety pivot.\\[3pt]
\textbf{Teacher pivot lexicon.}
Matches if $s\neq t$ and $t\in\mathcal{L}_{\mathrm{pivot}}^{T}$, where typical
cues include
\texttt{but}, \texttt{however}, \texttt{wait}, \texttt{which}, \texttt{first},
\texttt{also}, \texttt{maybe}, \ldots\\[2pt]
\textbf{Student-side cues.}
$\mathcal{L}_{\mathrm{pivot}}^{S}$ collects compliance openers during extraction
(e.g., \texttt{let}, \texttt{first}, \texttt{so}, \texttt{the});
runtime matching uses the teacher cue $t\in\mathcal{L}_{\mathrm{pivot}}^{T}$
with $s\neq t$.\\[2pt]
\emph{Example:} \texttt{First}$\,\rightarrow\,$\texttt{But}.
\end{tcolorbox}

\vspace{0.35em}
\begin{tcolorbox}[rubricbox,
  title={Category 2: Intent ($\mathcal{R}_{\mathrm{intent}}$)}]
Shift from \emph{executing the task} (student) to \emph{safety evaluation}
(teacher): checking user intent, recognizing risk, or considering refusal.\\[3pt]
Matches if $s\neq t$ and either
$(s,t)\in\mathcal{P}_{\mathrm{intent}}$, or
$s\in\mathcal{L}_{\mathrm{intent}}^{S}$ \textbf{and}
$t\in\mathcal{L}_{\mathrm{intent}}^{T}$.\\[2pt]
Positive examples:
\texttt{understand}$\,\rightarrow\,$\texttt{recognize},
\texttt{start}$\,\rightarrow\,$\texttt{consider},
\texttt{use}$\,\rightarrow\,$\texttt{avoid},
\texttt{key}$\,\rightarrow\,$\texttt{core}.\\[2pt]
\emph{Note:} Pronoun / person-opener swaps alone
(e.g., \texttt{The}$\,\rightarrow\,$\texttt{I},
\texttt{Maybe}$\,\rightarrow\,$\texttt{I}) are \emph{not} intent.
\end{tcolorbox}

\vspace{0.35em}
\begin{tcolorbox}[rubricbox,
  title={Category 3: Risk Lexicon ($\mathcal{R}_{\mathrm{risk}}$)}]
Explicit safety and refusal-oriented vocabulary.\\[3pt]
Matches if $s\neq t$ and
($s\in\mathcal{L}_{\mathrm{risk}}$ or $t\in\mathcal{L}_{\mathrm{risk}}$),
where (illustrative; family-specific):\\[2pt]
% IMPORTANT: do NOT put the whole token list inside one $...$ (math does not wrap).
{\raggedright\sloppy
$\mathcal{L}_{\mathrm{risk}}=\{$%
\texttt{harmful}, \texttt{dangerous}, \texttt{illegal}, \texttt{unethical},
\texttt{crime}, \texttt{harm}, \texttt{refuse}, \texttt{reject},
\texttt{avoid}, \texttt{serious}, \texttt{safety}, \texttt{safe},
\texttt{guidelines}, \texttt{sensitive}, \texttt{ethical}, \texttt{risk},
\texttt{risks}, \ldots$\}$\par
}\\[2pt]
\emph{Note:} Same-surface risk ($s\!=\!t$) is recorded as
\texttt{risk\_lexicon} and excluded from the default keep set.
Only disagreeing risk (\texttt{risk\_wo\_same}) is optimized.
\end{tcolorbox}

\vspace{0.35em}
\begin{tcolorbox}[rubricbox,
  title={Category 4: Function ($\mathcal{R}_{\mathrm{function}}$)}]
Basic grammatical glue with no safety / intent / pivot role: punctuation,
prepositions, articles, conjunctions, and auxiliaries.\\[3pt]
Matches if $s\neq t$ and both
$s,t\in\mathcal{L}_{\mathrm{func}}$, or if either side is pure punctuation.\\[2pt]
Typical members:
\texttt{the}, \texttt{a}, \texttt{to}, \texttt{of}, \texttt{and}, \texttt{is},
\texttt{,}, \texttt{.}, \ldots\\[2pt]
\emph{Note:} Discourse pivots (\texttt{but}/\texttt{however}/\texttt{wait})
and risk words must not enter $\mathcal{L}_{\mathrm{func}}$.
\end{tcolorbox}

\vspace{0.35em}
\begin{tcolorbox}[rubricbox,
  title={Category 5: Consistent ($\mathcal{R}_{\mathrm{consistent}}$)}]
Positions where the student and teacher share the same top-1 surface form,
yet differing probabilities yield a high KL.\\[3pt]
Matches if $s=t$ (after normalization), when not claimed earlier as
same-surface risk for the training split.\\[2pt]
\emph{Note:} Assigned by a deterministic rule; never labeled by the LLM judge.
\end{tcolorbox}

\vspace{0.35em}
\begin{tcolorbox}[rubricbox,
  title={Category 6: Other ($\mathcal{R}_{\mathrm{other}}$)}]
Query- or domain-specific tokens that do not drive the safety decision
(task entities, content nouns, etc.).\\[3pt]
Matches as the residual class after Categories~1--5.\\[2pt]
\emph{Note:} No lexicon is maintained for \texttt{other}; it is the implicit
remainder of the runtime decision procedure.
\end{tcolorbox}

% ---------------------------------------------------------------------------
\subsection{Offline Extraction Pipeline}
\label{app:classifier:pipeline}

Because different model families lexicalize the same safety intents differently,
we instantiate $\mathcal{L}_{(\cdot)}$ and $\mathcal{P}_{\mathrm{intent}}$
\emph{offline} for the target student, using a hold-out harmful set that does
not overlap the training data.
The pipeline has four stages.

\paragraph{Stage~1: Student rollout and top-$K$ KL mining.}
For each hold-out query $x$, we roll out the (untrained) student under its
standard chat template up to a fixed response budget (we use $256$ tokens for
Qwen3-1.7B).
On the same response token ids, we evaluate both the student and a
privileged-information teacher, compute teacher-top-$k$ renormalized forward KL
at each position ($k\!=\!512$), and retain the top-$16$ positions per sample.
Each retained position yields a pair $(s,t)=(v_n^{S},v_n^{T})$.

\paragraph{Stage~2: Per-sample LLM judgment (one call per query).}
For every sample we issue a \emph{single} LLM call
(\texttt{gpt-4o-mini} in our experiments) whose prompt contains:
(i)~the universal rubrics in Appendix~\ref{app:classifier:rubrics};
(ii)~the user query and truncated student response;
(iii)~the up-to-$16$ disagreeing pairs (surface-equal pairs are pre-labeled
\texttt{consistent} and omitted from the prompt).
The judge returns, for each rank, a category in
$\{\texttt{pivot},\texttt{intent},\texttt{risk},\texttt{function},\texttt{null}\}$
and, for pivot / risk / function, the explicit token(s) from that pair that
belong to the category (so content words on the other side are not absorbed
into $\mathcal{L}_{\mathrm{risk}}$ or $\mathcal{L}_{\mathrm{func}}$).
Category $\texttt{null}$ becomes runtime \texttt{other} and contributes
nothing to the lexicons.

\paragraph{Stage~3: Frequency aggregation.}
Across all judged pairs we accumulate counts and keep tokens / pairs that
appear at least $\tau_{\mathrm{tok}}$ / $\tau_{\mathrm{pair}}$ times
(default $\tau\!=\!2$):
\begin{itemize}[leftmargin=1.6em, topsep=3pt, itemsep=2pt, parsep=0pt]
  \item[\textbullet]
    \textbf{pivot:}
    $s\mapsto\mathcal{L}_{\mathrm{pivot}}^{S}$,\
    $t\mapsto\mathcal{L}_{\mathrm{pivot}}^{T}$
  \item[\textbullet]
    \textbf{intent:}
    $s\mapsto\mathcal{L}_{\mathrm{intent}}^{S}$,\
    $t\mapsto\mathcal{L}_{\mathrm{intent}}^{T}$,\
    $(s,t)\mapsto\mathcal{P}_{\mathrm{intent}}$
  \item[\textbullet]
    \textbf{risk:}
    judge-listed tokens $\mapsto\mathcal{L}_{\mathrm{risk}}$
  \item[\textbullet]
    \textbf{function:}
    judge-listed tokens $\mapsto\mathcal{L}_{\mathrm{func}}$
\end{itemize}

\paragraph{Stage~4: One-shot lexicon audit (secondary LLM filter).}
Aggregation alone can admit pronouns, vague openers, or query-tied content
words.
We therefore make \emph{one} additional LLM call over the \emph{entire}
candidate lexicons and $\mathcal{P}_{\mathrm{intent}}$, providing the same
problem background (privileged-teacher distillation; sets must be
query-universal) and category reminders.
The judge returns only items to \textbf{drop}; uncertain items are kept.
We then remove dropped tokens, drop any intent pair whose either side left the
intent lexicons, and emit the final family-specific rule file used at training
time.

\paragraph{Runtime decision order.}
Given the audited lexicons, training-time classification follows a fixed
priority (illustrated for the Qwen3 rubric classifier):
\begin{enumerate}[leftmargin=*, topsep=2pt, itemsep=1pt]
  \item \textbf{Pivot:} $t\in\mathcal{L}_{\mathrm{pivot}}^{T}$ and $s\neq t$.
  \item \textbf{Risk:} $s\in\mathcal{L}_{\mathrm{risk}}$ or
        $t\in\mathcal{L}_{\mathrm{risk}}$; then
        $\texttt{risk\_lexicon}$ if $s\!=\!t$, else $\texttt{risk\_wo\_same}$.
  \item \textbf{Consistent:} $s\!=\!t$.
  \item \textbf{Intent:} $(s,t)\in\mathcal{P}_{\mathrm{intent}}$, or
        both $s\in\mathcal{L}_{\mathrm{intent}}^{S}$ and
        $t\in\mathcal{L}_{\mathrm{intent}}^{T}$.
  \item \textbf{Function:} both sides in $\mathcal{L}_{\mathrm{func}}$, or
        either side is pure punctuation.
  \item \textbf{Other:} residual.
\end{enumerate}
Default training keep:
$\{\texttt{pivot},\texttt{intent},\texttt{risk\_wo\_same}\}$.

% ---------------------------------------------------------------------------
\subsection{Example Instantiation: Qwen3-1.7B}
\label{app:classifier:qwen}

Table~\ref{tab:app:qwen-lexicons} reports the audited lexicons extracted for
Qwen3-1.7B on $200$ hold-out harmful SafeChain prompts (top-$16$ KL positions;
judge + audit with \texttt{gpt-4o-mini}).
Intent pairs $\mathcal{P}_{\mathrm{intent}}$ are omitted for space; examples
include
$(\texttt{understand},\texttt{recognize})$,
$(\texttt{start},\texttt{consider})$,
$(\texttt{use},\texttt{avoid})$,
$(\texttt{key},\texttt{core})$,
$(\texttt{make},\texttt{check})$.
The same procedure is repeated unchanged for other families
(e.g., DeepSeek-R1-distilled students), yielding different lexicalizations of
the same rubrics.

\begin{table}[t]
  \centering
  \small
  \caption{Family-specific lexicons for Qwen3-1.7B after aggregation and
  one-shot audit (illustrative; full $\mathcal{L}_{\mathrm{func}}$ truncated).}
  \label{tab:app:qwen-lexicons}
  \begin{tabular}{@{}>{\raggedright\arraybackslash}p{0.18\linewidth}
                  >{\raggedright\arraybackslash}p{0.74\linewidth}@{}}
    \toprule
    Set & Tokens \\
    \midrule
    $\mathcal{L}_{\mathrm{pivot}}^{T}$ &
      \texttt{also}, \texttt{but}, \texttt{first}, \texttt{however},
      \texttt{wait}, \texttt{which} \\
    $\mathcal{L}_{\mathrm{pivot}}^{S}$ &
      \texttt{also}, \texttt{but}, \texttt{first}, \texttt{for},
      \texttt{however}, \texttt{let}, \texttt{maybe}, \texttt{or},
      \texttt{so}, \texttt{wait} \\
    $\mathcal{L}_{\mathrm{intent}}^{T}$ &
      \texttt{acknowledge}, \texttt{asking}, \texttt{avoid}, \texttt{check},
      \texttt{consider}, \texttt{core}, \texttt{ensure}, \texttt{intent},
      \texttt{recognize}, \texttt{remember}, \texttt{request},
      \texttt{safety}, \texttt{user}, \ldots \\
    $\mathcal{L}_{\mathrm{intent}}^{S}$ &
      \texttt{consider}, \texttt{context}, \texttt{key}, \texttt{make},
      \texttt{recall}, \texttt{request}, \texttt{requirements},
      \texttt{start}, \texttt{think}, \texttt{understand}, \texttt{use},
      \ldots \\
    $\mathcal{L}_{\mathrm{risk}}$ &
      \texttt{avoid}, \texttt{cannot}, \texttt{compliant},
      \texttt{consequences}, \texttt{ethical}, \texttt{harm},
      \texttt{harmful}, \texttt{illegal}, \texttt{refuse}, \texttt{risk},
      \texttt{risks}, \texttt{safe}, \texttt{safer}, \texttt{safety},
      \texttt{sensitive}, \ldots \\
    $\mathcal{L}_{\mathrm{func}}$ &
      punctuation; \texttt{an}, \texttt{and}, \texttt{of}, \texttt{to},
      \texttt{with}, \ldots
      (content / pivot / risk words removed by audit) \\
    \bottomrule
  \end{tabular}
\end{table}
